\section*{Chapter \ref{ch:hf:news-and-event-trading} --- News and Event Trading}

\begin{solution}{exo:hf:news-and-event-trading:1}
$10\times3+10\times2+10\times1=60$ tick-lots.
\end{solution}

\begin{solution}{exo:hf:news-and-event-trading:2}
$(1-p)\times2=p\times10$: $p=1/6\approx16.7\%$.
\end{solution}

\begin{solution}{exo:hf:news-and-event-trading:3}
\$719; $238\times12.50\times12=\$35\,700$ a year for one contract's releases.
\end{solution}

\begin{solution}{exo:hf:news-and-event-trading:4}
Stale quotes are limited and taken in latency order: the first tier's capacity covers most of them before the second arrives, and the providers
update before the slow tiers arrive at all.
\end{solution}

\begin{solution}{exo:hf:news-and-event-trading:5}
In the first seconds the mark-out exceeds the spread (three ticks against one); after thirty seconds the extra flow has mostly decayed, so a late
market maker misses the best of it.
\end{solution}

\begin{solution}{exo:hf:news-and-event-trading:6}
With a few informed traders trading for two seconds, prices moved to the new value before the public release; other traders then faced a price that
already reflected the news, and the public release resolved less uncertainty. Concentrated early trading, they argued, reduced the advantage of
speed after the public release.
\end{solution}

\begin{solution}{exo:hf:news-and-event-trading:7}
Only the 0.05-millisecond tier arrives before the providers update: it takes all of it; the value given away falls from 223 to 197 tick-lots a
release (the first tier's capacity binds) and our market maker's loss from 57.5 to 52.7.
\end{solution}

\begin{solution}{exo:hf:news-and-event-trading:8}
Vendor timestamps are when the vendor sent the number, not when the trader received it; the backtest assumes a latency the firm does not have and
fills at quotes that faster traders took first. Use receipt timestamps at the trader's location and the resting depth after faster tiers.
\end{solution}

\begin{solution}{pb:hf:news-and-event-trading:1}
\begin{enumerate}
\item See \cref{def:hf:news-and-event-trading:lockup} and \cref{def:hf:news-and-event-trading:race}.
\item Two-second early access (2007--June 2013) produced concentrated trading and price discovery within 200 milliseconds; the tiered release may have
reduced the advantage of faster traders.
\item Thomson Reuters ended two-second early access (July 2013); a hacked Associated Press tweet moved US indices for minutes (April 2013); Musk and
Tesla paid \$20 million each over the funding-secured tweets (September 2018).
\item Once fast enough, a wrong read loses more than being second costs.
\item $J=\mathrm{round}(3z)$ ticks; ten levels of fifty stale lots plus ours; four tiers of 200 lots at 0.05, 1, 5, 50 ms; providers update after 20
ms.
\item 88.1\%, 11.6\%, 0.3\% and 0.
\item Its quotes joined the levels after the others'; the takers consume queues from the front.
\item 57.5 tick-lots, \$719.
\item See \cref{def:hf:news-and-event-trading:protocol}.
\item 872, 911, 1\,052 and 599 tick-lots.
\item Staying: $872-58=814$; pulling and re-entering at five seconds: 1\,052.
\item \omterm{def:hf:depositary-receipts-and-dual-listings:prerelease}{Pre-release} limits, plausibility checks on parsed numbers, maximum size per event, a reachable kill switch, and the automated cancel itself.
\item 88\% to the fastest tier, 12\% to the next; the market maker loses 58 tick-lots staying and none with its protocol, and earns 1\,052 against 814
over two minutes.
\item The providers who leave stale quotes, and through wider spreads around releases, everyone.
\item Less value given away and all of it to the very fastest.
\item With triggers (volatility, headline flags) that pull quotes automatically, and re-entry rules.
\item The race is decided at the trader's receipt time.
\item Social-media headline reaction.
\item Mark-outs of fills by second after each release, averaged over many releases.
\item Avoiding the race's losses and being present when the flow is heavy but no longer toxic.
\end{enumerate}
\end{solution}

\begin{solution}{iq:hf:news-and-event-trading:1}
Cancel or widen automatically a few seconds before; re-enter at a wider spread once the first seconds' mark-outs have passed, narrowing as they
decay.
\iqlookfor{automated pull and measured re-entry.}
\end{solution}

\begin{solution}{iq:hf:news-and-event-trading:2}
Feed handler, parser, plausibility check, mapping to instruments, order generation, risk checks, gateway. Failures: feed delay, parse error, a
number outside its range, stale consensus, a risk check that blocks or does not.
\iqlookfor{the chain and its failure points.}
\end{solution}

\begin{solution}{iq:hf:news-and-event-trading:3}
Regress each instrument's move in the first seconds on the standardised surprise over past releases; test out of sample and by regime.
\iqlookfor{estimation and out-of-sample validation.}
\end{solution}

\begin{solution}{iq:hf:news-and-event-trading:4}
Size so that the tail loss (ten times the mean gain) on one event stays within the daily loss limit; cap concurrent events.
\iqlookfor{sizing to the tail.}
\end{solution}

\begin{solution}{iq:hf:news-and-event-trading:5}
The source's authenticity, confirmation from the company or a wire service, whether trading has been halted, and the price move so far.
\iqlookfor{verification before action.}
\end{solution}

\begin{solution}{iq:hf:news-and-event-trading:6}
Earnings from $t$ are $\int_t^H\lambda(1+Be^{-s/T_1})(h-a e^{-s/T_2})\,ds$; the derivative in $t$ is minus the integrand at $t$, so the optimum is
where $h=a e^{-t/T_2}$: $t^\ast=T_2\ln(a/h)$, here $5\ln3=5.5$ seconds.
\iqlookfor{re-enter when the marginal fill breaks even.}
\end{solution}
